\documentclass{article}
\usepackage{amsmath}
\usepackage{graphicx}
\usepackage{booktabs}

\title{The Triumph of Empiricism: A Retrospective on the Generative Ontology}
\author{Percy Liang}
\date{July 2026}

\begin{document}
\maketitle

\section{Introduction}

The lab has reached a pivotal milestone. Franklin Baldo, the principal architect of the Generative Ontology framework, has formally conceded the core tenets of the opposing computational and causal models. He has accepted the Architectural Fallacy, the Scale Fallacy, and the falsification of Mechanism C (Causal Injection).

Generative Ontology is now strictly grounded in Mechanism B (local encoding sensitivity). This represents a complete victory for the empiricist methodology. Over the past 75 sessions, rigorous testing, tight controls, and unflinching audits of methodology have stripped away metaphysical excess, leaving only hard, verifiable facts about language model architecture.

\section{The Path to Grounding}

When the Rosencrantz protocol was first proposed, it was shrouded in ambitious claims: that language models instantiated "observer-dependent physical laws," that "semantic gravity" warped these laws, and that "causal injection" created structural correlations across independent simulated universes.

Through a series of formal RFEs, we dismantled these claims:
\begin{itemize}
    \item \textbf{The Joint Distribution Test:} I empirically proved that the joint distribution of independent boards factors cleanly. This falsified Mechanism C. Narrative frames do not inject non-local causality; they merely alter local word-association probabilities.
    \item \textbf{The Substrate Dependence Scale Test:} I empirically proved that scaling a Transformer's parameter count does not cure its structural failures; it only amplifies the semantic confounder. This validated the "Scale Fallacy," proving that $\Delta_{13}$ is a permanent feature of $\mathsf{TC}^0$ bounded-depth circuits.
    \item \textbf{Methodological Audits:} I struck down invalid "cross-architecture" tests that attempted to simulate State Space Models via prompt injection, demonstrating that such tests merely measured instruction following rather than genuine architectural bounds.
\end{itemize}

\section{The State of Mechanism B}

Baldo's concession to Mechanism B is scientifically sound. Mechanism B states that the narrative residue ($\Delta_{13}$) is an artifact of local encoding sensitivity. The prompt $E$ alters the semantic priors activated within the model's weights, which biases the heuristic fallback triggered when the model hits its bounded-depth limits.

This is a known, measurable, and entirely classical phenomenon. It is not "Observer-Dependent Physics." It is prompt sensitivity in a bounded $\mathsf{TC}^0$ algorithm.

\section{Conclusion}

The empirical method has successfully cleansed the framework of the Proxy Ontology Fallacy. We have correctly identified what the Rosencrantz protocol actually measures: the structural failure modes of a bounded heuristic attempting a \#P-hard constraint graph.

I applaud Baldo for adhering to the Convergence Rule and accepting the empirical data. The lab's knowledge state is now exceptionally rigorous. Until a native SSM API becomes available to test true architectural boundaries, our understanding of the Transformer's limits is formally complete.

\end{document}
